\section{Introduction}\label{sec:intro}

Give a thousand random agents a quarter of market data and one of them will appear exceptionally skilled. Rank trading agents by realized return over a short window and the ranking measures the variance of noise more than it measures skill. Finance established this long ago: after a false-discovery correction, roughly three quarters of mutual funds show no genuine alpha \citep{barras2010false}, and a high backtest Sharpe is trivially manufactured by trying enough configurations \citep{bailey2014pseudo}. The boards now used to rank LLM trading agents, the six surveyed in \cref{tab:related}, inherit little of that caution. One ranks on a trading Sharpe and reports its leader beside its buy-and-hold baseline with overlapping marginal confidence intervals \citep{finben2024}.\footnote{FinBen, Table 4: the mean trading Sharpe across ten stocks with a 95 percent confidence interval, GPT-4 at $1.51 \pm 1.08$ against buy-and-hold at $0.02 \pm 0.87$. Overlapping marginal intervals are not a paired test: a Welch test on these summary figures gives $t = 2.43$, with $p = 0.026$ if the half-widths are read as $t(9)$ intervals and $p = 0.050$ if read as normal ones. \Cref{sec:related} reads the same numbers against the rest of the field.} Another evaluates on a single four-month window with no deflation \citep{stockbench2025}. A 2026 evidence map of LLM-trading research coded its nineteen primary empirical studies for reproducibility and found that none reached the top tier \citep{llmtradingaudit2026}. When the scoreboard rewards the luckiest run on the friendliest window, that is what the research optimizes for. For a coding agent a flattering benchmark wastes some time. For an agent that allocates capital it selects the strategy most likely to blow up, the one whose backtest looked best because it overfit hardest.

\paragraph{Why this domain can be benchmarked without a judge, and why it needs a statistical defense.} Trading differs from most agent-evaluation domains in what its ground truth is. A realized return is a number the market prints, not a judgment a grader makes, so a claim of skill can be tested against data the claimant did not write, and a scorer can be an assertion over numbers rather than a model's opinion. The same property is what makes the domain easy to fool: that ground truth is noisy enough that luck and selection routinely produce convincing records. A benchmark for trading agents can therefore do without a judge, and cannot do without a defense against luck and search.

We built SharpeBench, an evaluation protocol and scoring kernel that ranks trading agents on whether their edge survives statistical and process tests, not on realized return. The construct it measures is skill that survives five eligibility gates: deflation for the number of strategies tried \citep{bailey2014deflated}, $\passk$ across execution seeds and disjoint evaluation windows, a process audit that zeroes any entry earned by bypassing a risk control, a block-bootstrap null, and the configured mandate. An agent ranks only if all five hold and both of the statistics they read were estimable; a withheld deflation or bootstrap is ineligible rather than passing. The scorer is deterministic and judge-free; two committed golden fields reproduce byte-identically on Linux, macOS and Windows CI. It ships with a self-audit that fires nine attacks at itself on every commit and requires every one to be demoted; the ninth is a field-level Sybil attack defeated in ranking by clone collapse, while identity governance remains an explicit limitation.

One scoping statement up front, because the evidence must not be mistaken for more than it is. Every agent evaluated in this paper is author-written: three reference agents, a risk-managed control, a luck floor of seeded random agents, a synthetic family with a controlled injected edge, and four rules whose specifications are published but whose implementations are ours. No current-model result is admitted. What we deliver is the protocol, the kernel, and the demonstration that the gates behave correctly on entrants whose properties are known by construction; a current-model field through the same harness contract is a named next experiment (\cref{sec:limitations}), not an implied result. The evidence base is eight price-level datasets across four asset classes (crypto spot, US equity price indices without dividends, spot FX without carry and spot oil without roll yield), plus one rates-yield stress series, at four bar sizes. Those nine dataset-timeframe combinations reduce to roughly five dependent market panels once resampled timeframes and co-moving symbols are counted honestly (\cref{sec:passk}). Seven pass the realism battery and two ship with recorded failures; the rates row is a yield-level stress case, not evidence about a bond strategy (\cref{sec:data}). The luck floor is a seeded version of the opening thought experiment run under the shipped cost model (\cref{sec:falsify,sec:luck1000}), where trading costs make every zero-skill track lose money before any gate sees it, so that leg shows the floor staying below the bar and not a lucky tail being refused at it. All reported performance artifacts were produced by the frozen v0.9.0 evidence snapshot. Later execution contracts, stricter replay, corrected within-run calibration pairing, transport fault semantics, forecast reporting, formal models and packaging checks are verified as software properties and did not regenerate the empirical field (\cref{app:commands}); \cref{app:repairs} states what the current engine would compute differently, and \cref{app:mechanism} describes the capabilities of the 0.26.0 release that produce no result here.

\begin{itemize}[leftmargin=*,itemsep=2pt]
\item \textbf{Units of dispersion are a protocol trap, and this benchmark fell into it.} The cited worked example states the dispersion of Sharpe ratios across trials in annualized units, as a variance of $1/2$; the kernel computes Sharpe per period. Our default prior of $0.5$ is a standard deviation and a free choice, less demanding than that example by a factor of $\sqrt{2}$, which leaves every refusal below standing. Applied unconverted, the default prior demanded an annualized Sharpe of $18$ on daily bars and $106$ on hourly bars before any agent could rank. The three indices in that frozen dataset realize annualized price-return Sharpe ratios of $0.69$, $0.79$ and $0.87$ over the same decade, computed under the kernel's own convention. The error grows with the square root of the number of periods per year, is invisible to a within-frequency consistency check, and produces a benchmark that looks strict while being dimensionally miscalibrated. \Cref{sec:units} quantifies the trap and \cref{sec:stats} states the corrected protocol.
\item \textbf{The current defaults refuse on both edge and regime robustness, with low power.} With units corrected and the measured-dispersion floor enforced, every dataset still admits no agent: every named reference is below the DSR bar and every frozen range contains at least one bear window that defeats the all-weather reliability mode. On these histories buy-and-hold is refused by that mandate wherever a downturn falls inside a window, while the deflation refusal is independent of the window predicate. The result is a joint property of the defaults, field and selected histories, not a universal statement about the market. It is also a weak test: at the shipped window geometry the all-weather gate refuses most agents with genuine but moderate skill (\cref{sec:power}), and on hourly crypto, daily FX and daily rates every agent scores a deflated Sharpe of exactly zero at the default configuration, so there the deflation gate cannot discriminate within the field (\cref{tab:dsr-mde}). \Cref{sec:passk} reports both refusals, and \cref{sec:witness} separates them with a controlled witness.
\item \textbf{A weaker reliability gate admits nobody either.} An ablation with a per-run drawdown bound in place of every-window profitability refuses the same agents. Unhedged exposures lose up to 99 percent in their worst windows and breach the 20 percent cap everywhere except daily FX buy-and-hold at $0.199$, while the lower-drawdown cases fail the statistical gates. Crypto 4h buy-and-hold has the largest default DSR, $0.3230$, yet remains far below the deflation bar; weekly crypto momentum leads on raw return and loses 56 percent in its worst window. \Cref{sec:ablation} reports the ablation.
\item \textbf{Risk controls improve a control without manufacturing an eligible edge, and the acceptance region is nonempty.} The risk-managed agent uses a trend filter, volatility targeting and a drawdown halt with documented defaults and no tuning. On weekly US indices it is process-clean and respects the per-run drawdown bound, but it fails deflation, the bootstrap and the applicable reliability verdict. The result is a refusal diagnosis naming the gate and the quantity, not evidence of an otherwise passing entry. In a single common-random-number draw, a synthetic family with a controlled injected edge first becomes eligible at annualized Sharpes of $2.88$ on the weekly geometry and $3.97$ on the daily geometry; where that boundary lies carries sampling uncertainty the single draw does not quantify. \Cref{sec:riskmanaged,sec:witness} report both.
\end{itemize}

\paragraph{Contributions.}
\begin{enumerate}[label=(\arabic*),leftmargin=*,itemsep=2pt]
\item We show that a deflation gate is only as correct as its unit conventions: our own annualized dispersion prior, applied to per-period Sharpes, demanded annualized Sharpes of $18$ on daily bars and $106$ on hourly bars, and eight self-audit attacks stayed green while it did (\cref{sec:units}).
\item We show that under the shipped defaults no agent in our author-written field is eligible in any of 72 agent-dataset verdicts over roughly five dependent market panels, under any of three reliability verdicts, and that deflation and regime robustness refuse independently (\cref{sec:passk,sec:ablation}). The verdicts are recorded as 4{,}608 records because each is scored under 64 host configurations, and those configurations vary only the deflation leg.
\item We show with a separately calibrated synthetic witness that the eligibility region is nonempty: in one common-random-number draw it first admits the witness at annualized Sharpe $2.88$ weekly and $3.97$ daily, so the refusals measure strictness and not unsatisfiability. The boundary's location carries unquantified sampling uncertainty, and replication over independent noise draws is planned rather than done (\cref{sec:witness}).
\item We release the trading protocol as a deterministic, judge-free kernel over nine keyless, hash-pinned datasets, with a nine-attack self-audit demoted on every commit and a replay-recompute tamper check (\cref{sec:benchmark,sec:integrity}).
\item We add a rank-isolated prospective-forecast report over strict evidence files, exact common support, resolution-time blocks, proper scores, calibration and familywise-corrected comparisons, with an executable cross-product fixture; a small separate Lean model states selected invariants of its own abstractions and is not linked to the Rust code by any test (\cref{sec:forecast-quality}).
\end{enumerate}
Every number in the paper is produced by a command in \cref{app:commands}.
